% ============================================================
% SS-6: Deuteron Observables Beyond Binding:
%       Scope and Limits of the Base-to-Base Picture
% Strong Sector Series
%
% Version 0.4 --- 17 August 2026 (Patch 3213):
%   extended the generated identifier appendix to all five identifier
%   families (THEO-, FI-, PRED-, CONJ-, PH- as well as OPEN-), and
%   replaced review-convergence scores with AI review panel wording
%   where present. No claim, equation, number or section changed.
% Version 0.3 --- 17 August 2026 (Patch 3212):
%   added the generated plain-language appendix glossing the internal
%   programme identifiers used in this paper (founder jargon ruling, 16
%   Aug 2026). No claim, equation, number or section changed.
% Version 0.2 — 19 April 2026
%
% CHANGELOG:
%   v0.2 (19 Apr 2026) — AUTHOR-SIDE CORRECTNESS PASS before external review.
%     Three substantive fixes identified by Opus self-review:
%
%     FIX 1 (§4.2 Q_d comparison): v0.1 compared body-frame intrinsic
%     quadrupole Q_0^int = -0.22 fm^2 with lab-frame observed Q_d =
%     +0.286 fm^2 without converting frames. For a rigid prolate rotor
%     at J=K=1, Q_obs = Q_0 * J(2J-1)/[(J+1)(2J+3)] = Q_0/10. So if the
%     bipyramid dominated, it would predict Q_obs = -0.022 fm^2 --- an
%     order of magnitude below observation AND wrong sign. The
%     conceptual conclusion (Q_d is orbital-dominated) strengthens under
%     the correct framing. Section rewritten to present body-frame and
%     lab-frame quantities with appropriate conversion.
%
%     FIX 2 (§4.5 effective-range inversion): v0.1 used the wrong
%     algebraic inversion r_0 = 2(a - 1/kappa), giving r_0 = 2.21 fm vs
%     experimental 1.749 fm (claimed 26% error). The correct inversion
%     from 1/a = kappa - r_0 kappa^2/2 is r_0 = 2(kappa - 1/a)/kappa^2
%     = 1.76 fm --- a +0.8% match with experiment, not a 26% failure.
%     This flips the rhetorical effect: the effective-range expansion
%     works extremely well for the deuteron at leading order. Text and
%     framing adjusted accordingly.
%
%     FIX 3 (§2.3 kinetic energy): v0.1 said "kinetic energy ~20 MeV,
%     potential depth ~35 MeV" for the deuteron. Standard textbook
%     values are T ~ 33 MeV, V ~ 35 MeV (both large, nearly cancelling
%     to give 2.2 MeV net binding). Minor fix; no conceptual impact.
%
%     FRAMING SHARPENING: The Bethe-Peierls a_np = 1/kappa = 4.32 fm
%     result is universal two-body scattering theory, not a CPP
%     prediction. v0.1 was careful but could be read as claiming CPP
%     credit for a textbook result. v0.2 reframes this as "CPP supplies
%     B_d; two-body scattering theory supplies a_np via Bethe-Peierls."
%     What CPP adds (via OPEN-SS-20) is V_SR(r) and hence r_0 and a_np
%     beyond zero range.
%
%     NO substantive changes to: three-category classification,
%     OPEN-SS-20/21 registration, PROP-SS-6-1/2, abstract conclusions.
%     The paper's conceptual structure is unchanged; only specific
%     numerical claims and framing are tightened.
%
%   v0.1 (18 Apr 2026) — INITIAL SCOPING DRAFT. Follows SS-5 v5/v6's
%     success on light-nuclei binding by testing the SAME base-to-base
%     bipyramid geometry against additional deuteron observables.
%     SCOPE DECISION: After numerical exploration, the bipyramid cleanly
%     predicts (a) binding (done in SS-5), (b) zero-range scattering
%     length a_np = 1/kappa from B_d input alone (via standard Bethe-
%     Peierls relation, no CPP-specific derivation).
%     NEGATIVE FINDING: rigid bipyramid's intrinsic quadrupole moment
%     has the WRONG SIGN (oblate, -0.22 fm^2) vs observed prolate Q_d =
%     +0.286 fm^2. This reveals Q_d is ORBITAL-DOMINATED, not bipyramid-
%     dominated. Same conclusion transfers to r_d, P_D, mu_d — all 
%     require the deuteron orbital wavefunction at ~2 fm scale, which
%     CPP has not yet derived. 
%     Paper is a HONEST SCOPING REPORT that classifies deuteron 
%     observables into (A) bipyramid-geometric, (B) bipyramid-via-V_SR,
%     (C) orbital-dominated, and registers open problems for each.
%     Registers: OPEN-SS-20 (V_SR(r) shape derivation),
%                OPEN-SS-21 (deuteron orbital wavefunction from CPP),
%                PROP-SS-6-1 (Q_d orbital-dominated),
%                PROP-SS-6-2 (zero-range a_np relation).
%
% CONTRIBUTORS:
%   Thomas Lee Abshier ND — direction (bipyramid as multi-observable
%     test of base-to-base), scope authority (17-18 April 2026).
%   Claude Opus (Anthropic) — numerical exploration, Q_d sign finding,
%     scope recommendation, v0.1 drafting (18 April 2026).
%
% Series: Strong Sector
% Dependencies: SS-5 v6 (base-to-base bipyramid, B_d LO, cascade
%   formula, epsilon_d correction program), SS-2 (nucleon structure,
%   l_unit, l_edge, cage geometry), SM-8 (M_0 quantum).
% Resolves: nothing fully; partially addresses OPEN-SS-10 
%   (nuclear potential V(r) shape at short range).
% Registers: OPEN-SS-20 (V_SR shape), OPEN-SS-21 (orbital wavefunction),
%   PROP-SS-6-1 (Q_d orbital-dominated), PROP-SS-6-2 (zero-range a_np).
% ============================================================

\documentclass[11pt,letterpaper]{article}
\usepackage[utf8]{inputenc}
\usepackage[margin=1in]{geometry}
\usepackage[T1]{fontenc}
\usepackage{amsmath,amssymb,amsfonts,amsthm}
\usepackage{graphicx}
\usepackage{xcolor}
\usepackage{hyperref}
\usepackage{booktabs}
\usepackage{array}
\usepackage{enumitem}
\usepackage{float}
\usepackage[font=small, labelfont=bf]{caption}
\usepackage{parskip}
\usepackage{mdframed}

\hypersetup{
    colorlinks=true,
    linkcolor=blue,
    citecolor=blue,
    urlcolor=blue
}

\newtheorem{theorem}{Theorem}[section]
\newtheorem{proposition}[theorem]{Proposition}
\newtheorem{lemma}[theorem]{Lemma}
\newtheorem{corollary}[theorem]{Corollary}
\newtheorem{conjecture}[theorem]{Conjecture}
\newtheorem{definition}[theorem]{Definition}
\newtheorem{remark}[theorem]{Remark}
\newtheorem{openproblem}{Open Problem}
\newtheorem{finding}[theorem]{Finding}

\newcommand{\phig}{\varphi}
\newcommand{\Mzero}{M_0}
\newcommand{\Bpair}{B_{\mathrm{pair}}}
\newcommand{\Bd}{B_{d}}
\newcommand{\Qd}{Q_{d}}
\newcommand{\PD}{P_{D}}
\newcommand{\mud}{\mu_{d}}
\newcommand{\anp}{a_{np}}
\newcommand{\rzero}{r_{0}}
\newcommand{\lunit}{l_{\mathrm{unit}}}
\newcommand{\ledge}{l_{\mathrm{edge}}}
\newcommand{\VSR}{V_{\mathrm{SR}}}
\newcommand{\epsd}{\varepsilon_{d}}
\newcommand{\qCP}{q\mathrm{CP}}

\title{\textbf{SS-6: Deuteron Observables Beyond Binding:\\ Scope and Limits of the Base-to-Base Picture}\\[6pt]
{\large 600-Cell Standard Model Emergence Series}\\[4pt]
{\normalsize Version 0.2 --- 19 April 2026 (pre-review correctness pass)}}

\author{Thomas Lee Abshier, ND \and Claude Opus (Anthropic)}

\date{\normalsize Hyperphysics Institute, Kalispell, Montana\\
\url{https://hyperphysics.com}\\
\href{mailto:drthomas007@protonmail.com}{drthomas007@protonmail.com}\\[4pt]
\small OSF DOI: 10.17605/OSF.IO/JXE8D\\[2pt]
{\small Version~0.3 --- 17 August 2026 (Patch 3212: added the generated plain-language appendix glossing the internal programme identifiers used in this paper (founder jargon ruling, 16 Aug 2026). No claim, equation, number or section changed.)}\\[2pt]
{\small Version~0.4 --- 17 August 2026 (Patch 3213: extended the generated identifier appendix to all five identifier families (THEO-, FI-, PRED-, CONJ-, PH- as well as OPEN-), and replaced review-convergence scores with AI review panel wording where present. No claim, equation, number or section changed.)}}

\begin{document}
\maketitle

\begin{abstract}
SS-5 derived the deuteron binding energy $\Bd = 2.342$~MeV at leading order ($+5.3\%$) from base-to-base bonding of hybrid-tetrahedral nucleons, along with concurrent predictions for ${}^3$H, ${}^3$He, and ${}^4$He within $1.4\%$. A natural next question is whether the same bipyramid geometry predicts additional deuteron observables --- the quadrupole moment $\Qd$, D-state admixture $\PD$, magnetic moment $\mud$, matter radius $r_d$, and low-energy scattering parameters $(\anp, \rzero)$ --- with comparable zero-parameter economy.

This paper reports a scoping exploration with one definitive negative finding and two quantitative consistency checks. The negative finding: the rigid bipyramid's \emph{intrinsic} body-frame quadrupole is $Q_0^{\mathrm{body}} = -0.22$~fm$^2$ (oblate), which converts to a lab-frame spectroscopic contribution of only $\approx -0.022$~fm$^2$ --- an order of magnitude smaller than and opposite in sign to the observed prolate $\Qd = +0.286$~fm$^2$. This is not a failure of the mechanism; it reveals that $\Qd$, $r_d$, $\PD$, and $\mud$ are dominated by the \emph{orbital} deuteron wavefunction at $r \sim 2$~fm, not by the bipyramid at $\sim 0.4$~fm. The two consistency checks: (i) the universal Bethe-Peierls relation $\anp \approx 1/\kappa$ gives $\anp = 4.32$~fm from $\Bd$ alone ($-20\%$ vs observed 5.43~fm); this is standard two-body quantum mechanics, not a CPP-specific prediction. (ii) Inverting the effective-range expansion on the observed $\anp$ and $\kappa$ recovers the effective range to $+0.8\%$ ($\rzero = 1.76$~fm vs experimental 1.749~fm) --- a notable consistency check indicating the deuteron's short-range potential is smooth and well-characterized by a single length scale.

We classify each deuteron observable into three categories: (A) bipyramid-geometric and hence computable at SS-5 precision (binding, spin-parity, isospin); (B) bipyramid-via-$\VSR(r)$ and hence computable once the short-range potential shape is derived ($\anp$, $\rzero$, virtual singlet); (C) orbital-dominated and hence requiring the full deuteron wavefunction beyond the bipyramid core ($\Qd$, $r_d$, $\PD$, $\mud$). Two new open problems are registered: OPEN-SS-20 (derive $\VSR(r)$ from CPP primitives) and OPEN-SS-21 (derive the deuteron orbital wavefunction in the CPP framework). This paper does not attempt to close either.
\end{abstract}

\medskip\noindent
\textbf{Keywords:} deuteron quadrupole moment, scattering length, effective range, D-state admixture, magnetic moment, intrinsic vs orbital contributions, short-range potential, Bethe-Peierls relation, zero-range approximation, bipyramid geometry, base-to-base nucleon configuration, 600-cell lattice, Conscious Point Physics, scoping analysis.

\bigskip\noindent
\textbf{Plain Language Summary:}
SS-5 showed that gluing two tetrahedral nucleons together face-to-face correctly predicts their binding energy. A natural follow-up question is whether the same glued-tetrahedron geometry predicts \emph{other} deuteron properties --- its shape (quadrupole moment), its size (matter radius), and its scattering behaviour. This paper reports what we find when we ask this: the geometry predicts the binding cleanly, and gives a rough first-guess for how slowly-moving neutrons scatter off protons, but it cannot by itself predict the deuteron's shape and size, because the deuteron spends most of its time with the two nucleons orbiting \emph{around} each other at about 2~femtometres separation --- much larger than the half-femtometre bipyramid core. So this paper is an honest scoping report: here is what the geometric picture predicts well, here is what it doesn't, and here are the new pieces we would need to derive to get the rest.

\bigskip
\raggedright
\tableofcontents
\newpage

% ===================================================================
\section{Introduction}
\label{sec:intro}
% ===================================================================

SS-5 \cite{abshier2026ss5} derived four light-nuclei binding energies and three structural unboundness results from a single base-to-base bipyramid mechanism with zero fitted parameters. The question addressed here is whether the same mechanism predicts additional deuteron observables. If the bipyramid is a real geometric object and not a convenient mnemonic for the binding, then its charge distribution, spin structure, and coupling to external probes should produce definite predictions for $\Qd$, $\PD$, $\mud$, $r_d$, $\anp$, and $\rzero$ --- all precisely measured.

This question is motivated by the ``more stars for better navigation'' principle: a mechanism that predicts one observable cleanly but fails on others is partially wrong; a mechanism that predicts several observables concurrently within the same residual band is likely structural. SS-5's seven concurrent predictions (four quantitative bindings plus three unboundnesses) already constitute a strong test. Additional deuteron observables, if predicted at comparable precision from the same bipyramid, would further constrain the mechanism.

\subsection{The finding, briefly}

The result of this scoping is less triumphant than SS-5. Numerical exploration (§\ref{sec:observables}) produces two clear outputs:

\begin{enumerate}
\item \textbf{The rigid bipyramid cannot produce the observed $\Qd$.} The three net-$+1/3$ EM charges in the base-to-base contact plane give a body-frame intrinsic quadrupole $Q_0^{\mathrm{body}} = -0.22$~fm$^2$ (oblate), which converts for a $J=K=1$ rigid rotor to a lab-frame spectroscopic contribution of only $\approx -0.022$~fm$^2$. The observed $\Qd = +0.286$~fm$^2$ is ten times larger and opposite in sign. Interpreted correctly, this result \emph{reveals} that the observed $\Qd$ cannot arise from the bipyramid core; it must come from the \emph{orbital} wavefunction at the deuteron's ${\sim}2$~fm scale, where the proton and neutron spend most of their time separated rather than in contact.

\item \textbf{The Bethe-Peierls relation places $\Bd$ and $\anp$ in a universal consistency check.} For any shallow two-body bound state, $\anp \approx 1/\kappa$ with $\kappa = \sqrt{2\mu\Bd}/\hbar c$. Using $\Bd$ alone gives $\anp^{(0)} = 4.32$~fm against the measured 5.43~fm ($-20\%$). This is not a CPP-specific prediction; it is universal quantum scattering theory. What CPP \emph{has} supplied is $\Bd$; what it \emph{owes} (OPEN-SS-20) is $\VSR(r)$, from which $\rzero$ and the finite-range-corrected $\anp$ follow.

\item \textbf{The effective-range expansion works to 0.8\% at leading order.} Inverting the Bethe-Peierls relation on the observed $\anp = 5.425$~fm and $\kappa = 0.2316$~fm$^{-1}$ gives $\rzero = 2(\kappa - 1/\anp)/\kappa^2 = 1.76$~fm, against experimental $\rzero^{\mathrm{exp}} = 1.749$~fm. This $+0.8\%$ consistency is a notable positive finding: it indicates the deuteron's short-range potential is smooth and well-characterized by a single length scale $\rzero$, which whatever $\VSR(r)$ CPP eventually derives must reproduce.
\end{enumerate}

The upshot is a classification: some deuteron observables are bipyramid-computable; some require an additional CPP input (the $\VSR(r)$ shape); some require the orbital wavefunction and are beyond the present reach.

\subsection{What SS-6 delivers}

\begin{itemize}
\item A definitive structural finding: the observed deuteron quadrupole moment is orbital-dominated, not bipyramid-dominated. Rigid bipyramid alone predicts $\approx -0.022$~fm$^2$ (oblate) in the lab frame; observation is $+0.286$~fm$^2$ (prolate).
\item Two universal scattering-theory consistency checks: $\anp = 1/\kappa = 4.32$~fm from $\Bd$ alone ($-20\%$ vs measured 5.43~fm) and $\rzero = 1.76$~fm from effective-range inversion ($+0.8\%$ vs measured 1.749~fm).
\item A three-category classification of deuteron observables by derivability from the bipyramid.
\item Two new open problems: OPEN-SS-20 ($\VSR(r)$ shape) and OPEN-SS-21 (deuteron orbital wavefunction).
\end{itemize}

\subsection{What SS-6 does not deliver}

\begin{itemize}
\item No prediction of $\Qd$, $\PD$, $\mud$, or $r_d$ from CPP --- these are orbital-dominated.
\item No derivation of $\VSR(r)$ shape from CPP primitives.
\item No numerical computation of $\rzero$ or the effective-range-corrected $\anp$.
\item No improvement to SS-5's $\Bd$ prediction.
\end{itemize}

This is a scoping paper, not a result paper. Its purpose is to sharpen the programme's claims by distinguishing what the bipyramid mechanism actually predicts from what requires additional CPP input.

\subsection{Open problems addressed}

\begin{itemize}
\item \textbf{OPEN-SS-10} (Nuclear binding $V(r)$): partially clarified --- the integrated binding at A=2,3,4 was resolved in SS-5; the short-range shape of $\VSR(r)$ is now separately registered as OPEN-SS-20.
\item \textbf{New OPEN-SS-20}: derive the shape of $\VSR(r)$ as a function of inter-nucleon separation from CPP primitives. Controls $\anp$, $\rzero$, and the singlet virtual state.
\item \textbf{New OPEN-SS-21}: derive the deuteron orbital wavefunction at ${\sim}2$~fm from CPP in a framework that connects to the bipyramid core at ${\sim}0.4$~fm. Controls $\Qd$, $r_d$, $\PD$, $\mud$.
\item \textbf{New PROP-SS-6-1}: the rigid bipyramid's body-frame $Q_0^{\mathrm{body}} = -0.22$~fm$^2$ is oblate, lab-frame contribution $\approx -0.022$~fm$^2$; observed prolate $\Qd = +0.286$~fm$^2$ is orbital-dominated.
\item \textbf{New PROP-SS-6-2}: universal Bethe-Peierls gives $\anp = 1/\kappa = 4.32$~fm ($-20\%$) from $\Bd$; effective-range inversion additionally gives $\rzero = 1.76$~fm ($+0.8\%$).
\end{itemize}

% ===================================================================
\section{Setup: The Bipyramid from SS-5}
\label{sec:setup}
% ===================================================================

\subsection{Geometric summary}

In SS-5, the deuteron is the face-to-face contact of two hybrid-tetrahedral nucleons. Each nucleon (per SS-2) has:
\begin{itemize}
\item three quarks at the vertices of a triangular base face,
\item one ``open'' vertex opposite the base carrying net polarity ($+$ for proton, $-$ for neutron) but no additional electromagnetic charge,
\item isosceles base geometry: u-u edge stretched to $r_{uu} = (1 + \varepsilon_{\mathrm{cage}}) \ledge = 1.07$~fm with $\varepsilon_{\mathrm{cage}} = 1.94$; u-d edges at $r_{ud} = 0.62$~fm (SS-2 \cite{abshier2026ss2}).
\end{itemize}

In the base-to-base deuteron configuration, the proton base $\{u, u, d\}$ contacts the neutron base $\{d, d, u\}$ with the three charge-opposite vertex pairs aligned, forming a K$_3$ contact face. Three quark-quark DP chains span the contact, with K$_3$ collective mode reducing the three oscillators to one effective binding mode at $\Bpair = \Mzero/\phig = 2.342$~MeV (LO).

The open vertices of the two nucleons point \emph{outward} along the symmetry axis: proton's $+$ vertex at $z = +h$, neutron's $-$ vertex at $z = -h$, with $h$ the apex-to-base distance of a single cage. The resulting compound object is a \emph{bipyramid}: two tetrahedral cages joined at a common triangular base, apices pointing outward.

\subsection{Key geometric lengths}

For the numerical scale of the bipyramid core, SS-2 \cite{abshier2026ss2} and SM-8 \cite{abshier2026sm8} give:
\begin{align*}
\ledge &= 0.364 \text{ fm}, \\
\lunit &= \hbar c / \Lambda_{\mathrm{QCD}} = 0.589 \text{ fm}, \\
r_{uu} &= 1.07 \text{ fm} \quad \text{(proton u-u edge)}, \\
r_{ud} &= 0.62 \text{ fm} \quad \text{(u-d edges)}, \\
h &\approx 0.63 \text{ fm} \quad \text{(apex-to-base; equivalent-regular-tet estimate)}.
\end{align*}
The apex-to-apex long axis of the bipyramid is $2h \approx 1.26$~fm. The equatorial extent is set by $r_{uu} = 1.07$~fm. The bipyramid is mildly prolate, aspect ratio ${\sim}1.18$.

These scales are all ${\lesssim}1.3$~fm --- substantially smaller than the observed deuteron matter radius $r_d = 1.9753$~fm and charge radius $r_c = 2.1413$~fm.

\subsection{The separation of scales}

A crucial physical fact: in the deuteron, the proton and neutron are \emph{not} locked in the contact bipyramid geometry. They orbit at an rms separation of ${\sim}4$~fm (inferred from $r_d^2 = \langle r_{np}^2 \rangle/4$, giving $\langle r_{np}^2 \rangle^{1/2} \approx 3.95$~fm). This is $\sim 10\times$ the bipyramid's internal scale and $\sim 4\times$ the bipyramid's apex-to-apex dimension.

This separation of scales is not a CPP issue --- it is a standard feature of the deuteron recognized in all nuclear-physics treatments. The deuteron is a \emph{loosely-bound} state: $\Bd = 2.22$~MeV, kinetic energy $T \sim 33$~MeV, potential depth $V \sim 35$~MeV (standard textbook values), giving a binding $B = V - T \approx 2$~MeV that is the difference of two large, nearly-cancelling energies. The wavefunction extends far outside the attractive-potential range. The asymptotic behaviour $\psi(r) \sim e^{-\kappa r}/r$ with $1/\kappa = 4.32$~fm sets the observed size.

This separation of scales is the dominant physics controlling $\Qd$, $r_d$, $\PD$, and $\mud$. It is \emph{not} controlled by the bipyramid core.

% ===================================================================
\section{The Observables}
\label{sec:observables}
% ===================================================================

We work through the non-binding deuteron observables in turn, identifying what the bipyramid can and cannot say.

\subsection{Spin, isospin, parity}

Already addressed in SS-5 §7 at the bipyramid level: $J^P = 1^+$, $I = 0$. These are geometric consequences of the base-to-base K$_3$ alignment and the S-wave relative motion. No additional work required. \textbf{Category A (bipyramid-geometric).}

\subsection{Quadrupole moment $\Qd$}
\label{sec:Qd}

\subsubsection{Intrinsic bipyramid quadrupole (body frame)}

Consider the bipyramid at rigid contact ($s = 0$ gap between the two bases). The three quark-quark DP chains span the contact with opposite-charge pairs superposed at identical $(x, y)$ positions on the $z = 0$ contact plane:
\begin{itemize}
\item vertex 1: proton $u$ ($+2/3$) superposed with neutron $d$ ($-1/3$), net charge $+1/3$ at $(-r_{uu}/2, 0, 0)$;
\item vertex 2: proton $u$ ($+2/3$) superposed with neutron $d$ ($-1/3$), net charge $+1/3$ at $(+r_{uu}/2, 0, 0)$;
\item vertex 3: proton $d$ ($-1/3$) superposed with neutron $u$ ($+2/3$), net charge $+1/3$ at $(0, y_3, 0)$,
\end{itemize}
where $y_3 = \sqrt{r_{ud}^2 - (r_{uu}/2)^2} = 0.313$~fm is the base-triangle height from the u-u edge to the d vertex. Total net charge is $3 \times (+1/3) = +1$, the deuteron charge.

The apices at $z = \pm h$ carry polarity (one extra $\qCP$ each), not electromagnetic charge. They do not contribute to the intrinsic charge quadrupole.

In the \emph{body frame} (bipyramid symmetry axis aligned with $z$), with charges at positions $\vec r_i$ and values $q_i$, the intrinsic quadrupole moment is
\[
Q_0^{\mathrm{body}} = \sum_i q_i (3 z_i^2 - r_i^2).
\]
With all three charges at $z_i = 0$, the formula reduces to
\[
Q_0^{\mathrm{body}} = -\sum_i q_i (x_i^2 + y_i^2) = -\frac{1}{3}\left[2 \cdot \frac{r_{uu}^2}{4} + y_3^2\right] = -\frac{1}{3}\left[\frac{(1.07)^2}{2} + (0.313)^2\right] = -0.224 \text{ fm}^2.
\]

The negative sign indicates an \emph{oblate} intrinsic charge distribution: charge concentrated equatorially, neutral poles axially.

\subsubsection{Conversion to lab-frame spectroscopic quadrupole}

The observed deuteron quadrupole $\Qd = +0.286$~fm$^2$ \cite{pdg2024} is the \emph{spectroscopic} quadrupole in the $J=1, M=\pm 1$ stretched state. For a rigid rotor with symmetry axis $K = J$, the conversion from body-frame intrinsic $Q_0^{\mathrm{body}}$ to lab-frame spectroscopic $Q$ is standard \cite{krane_nuclear}:
\[
Q = Q_0^{\mathrm{body}} \cdot \frac{J(2J-1)}{(J+1)(2J+3)}.
\]
For $J = K = 1$: $Q = Q_0^{\mathrm{body}} / 10$. Applying this to the bipyramid's $Q_0^{\mathrm{body}} = -0.224$~fm$^2$ gives the spectroscopic value one would observe \emph{if the bipyramid dominated the deuteron's charge distribution}:
\[
Q^{\mathrm{bipyramid-only}} = \frac{-0.224}{10} \approx -0.022 \text{ fm}^2.
\]

\begin{finding}[Bipyramid cannot produce the observed $\Qd$]
\label{find:Qd_sign}
If the rigid base-to-base bipyramid dominated the deuteron's charge distribution, the predicted lab-frame quadrupole would be $Q^{\mathrm{bipyramid-only}} \approx -0.022$~fm$^2$: an order of magnitude smaller than the measured $\Qd = +0.286$~fm$^2$, \emph{and} the opposite sign. The bipyramid's equatorial-charge / neutral-pole geometry is intrinsically oblate; the observed deuteron charge distribution is prolate. This two-way mismatch cannot be closed by finite corrections to the bipyramid geometry.
\end{finding}

\subsubsection{Interpretation}

This is not a failure of the bipyramid mechanism. It is a diagnostic. The observed $\Qd$ is well-known in conventional nuclear physics to arise from the D-wave ($L = 2$) component of the deuteron wavefunction at orbital separation $r_{np} \sim 2$~fm. The relationship (standard \cite{krane_nuclear}):
\[
\Qd \approx \frac{\sqrt 2}{10} \int_0^\infty u(r) w(r) r^2 \, dr - \frac{1}{20}\int_0^\infty w(r)^2 r^2 \, dr,
\]
where $u(r)$ is the S-wave radial wavefunction and $w(r)$ is the D-wave wavefunction, weights $\Qd$ heavily with the large-$r$ structure of the wavefunction. With $\langle r^2 \rangle \approx 15$~fm$^2$ for the deuteron, the dominant contribution to $\Qd$ comes from $r \gtrsim 2$~fm, far outside the ${\lesssim}1$~fm bipyramid core.

The bipyramid's lab-frame spectroscopic contribution is $\approx -0.022$~fm$^2$ (oblate, small). The orbital wavefunction must therefore supply the entire observed $+0.286$~fm$^2$ plus a small compensating $+0.022$~fm$^2$ to cancel the bipyramid's intrinsic oblateness --- a total orbital contribution of $\approx +0.308$~fm$^2$. The bipyramid contribution is only ${\sim}7\%$ of the orbital contribution by magnitude; this is why the orbital regime dominates $\Qd$ observationally.

\begin{proposition}[PROP-SS-6-1: $\Qd$ is orbital-dominated]
\label{prop:Qd_orbital}
The rigid base-to-base bipyramid produces an oblate body-frame $Q_0^{\mathrm{body}} = -0.22$~fm$^2$, converting to a lab-frame spectroscopic contribution $\approx -0.022$~fm$^2$ --- an order of magnitude smaller than and opposite in sign to the observed prolate $\Qd = +0.286$~fm$^2$. The observed quadrupole therefore cannot arise from the bipyramid core alone; it is dominated by the orbital-scale ($r \gtrsim 2$~fm) D-wave wavefunction. Quantitative prediction of $\Qd$ from CPP requires OPEN-SS-21 (orbital wavefunction), not bipyramid geometry alone.
\end{proposition}

\textbf{Category C (orbital-dominated).}

\subsection{Matter and charge radii $r_d$, $r_c$}

$r_d = 1.975$~fm and $r_c = 2.141$~fm are both significantly larger than the bipyramid extent ($\lesssim 1.3$~fm). In standard nuclear physics, $r_d^2 = \frac{1}{4}\langle r_{np}^2 \rangle + \frac{1}{2}(r_p^2 + r_n^2)$, with the first term (relative-motion) dominating.

The relative-motion contribution is controlled by the asymptotic form $\psi(r) \propto e^{-\kappa r}/r$ with $1/\kappa \approx 4.32$~fm. The size of the deuteron is essentially the size of this exponential tail, not the bipyramid core. \textbf{Category C (orbital-dominated).}

\subsection{D-state admixture $\PD$ and magnetic moment $\mud$}

The D-state admixture is the fractional weight of the $L=2$ component of the deuteron wavefunction. Modern NN-potential extractions give $\PD = 4.5$--$5.8\%$ (Argonne v18 \cite{argonne_v18}: $5.76\%$; CD-Bonn: $4.85\%$; chiral N$^3$LO \cite{machleidt_entem}: $4.5$--$4.9\%$). The D-wave extends throughout the orbital region; its weight is fixed by the full wavefunction, not by the bipyramid's intrinsic angular structure.

The magnetic moment $\mud$ follows from $\PD$ via standard spin algebra:
\[
\mud = \mu_p + \mu_n - \frac{3}{2} \PD (\mu_p + \mu_n - \tfrac{1}{2}),
\]
with $\mu_p = 2.793 \mu_N$, $\mu_n = -1.913 \mu_N$ \cite{pdg2024}. The experimental $\mud = 0.8574 \mu_N$ implies $\PD \approx 3.9\%$ from this relation (slightly below the NN-potential extraction range, with the discrepancy attributed to meson-exchange and relativistic corrections).

Both $\PD$ and $\mud$ are downstream consequences of the orbital wavefunction. \textbf{Category C.}

\subsection{Scattering length $\anp$ and effective range $\rzero$}
\label{sec:anp_r0}

Low-energy np scattering in the triplet channel is described by the effective-range expansion
\[
k \cot \delta_0 = -\frac{1}{\anp} + \frac{1}{2} \rzero k^2 + \ldots
\]
For a loosely-bound two-body state, the Bethe-Peierls relation connects $\anp$ to the bound-state wavenumber $\kappa = \sqrt{2\mu\Bd}/\hbar c$:
\[
\frac{1}{\anp} = \kappa - \frac{1}{2}\rzero \kappa^2 + \ldots,
\]
so that
\[
\anp = \frac{1}{\kappa}\left(1 + \frac{\rzero \kappa}{2} + \ldots\right).
\]

With $\Bd = 2.2246$~MeV \cite{ame2020} and reduced mass $\mu = m_N/2 = 469.45$~MeV:
\[
\kappa = \frac{\sqrt{2 \mu \Bd}}{\hbar c} = 0.2316 \text{ fm}^{-1}, \quad \frac{1}{\kappa} = 4.318 \text{ fm}.
\]

\begin{proposition}[PROP-SS-6-2: Zero-range scattering length follows from $\Bd$ via universal two-body theory]
\label{prop:anp_zero_range}
For any shallow two-body bound state, the Bethe-Peierls zero-range limit gives
\[
\anp^{(0)} = \frac{1}{\kappa} = 4.318 \text{ fm}
\]
using $\Bd$ as the only input. This is not a CPP-specific prediction --- it is universal two-body quantum-mechanical scattering theory applied to any theory that correctly predicts $\Bd$. The measured value is $\anp = 5.425$~fm \cite{anp_reference}, giving a zero-range error of $-20.4\%$ that is the well-known finite-range shortfall of the deuteron.
\end{proposition}

\subsubsection{What CPP owes, and what CPP has already supplied}

CPP's path to a precision $\anp$ prediction is structural: CPP supplies $\Bd$ (done in SS-5); two-body scattering theory supplies $\anp^{(0)} = 1/\kappa$ automatically (universal); and CPP must additionally supply $\VSR(r)$ (OPEN-SS-20) to reach $\rzero$ and the finite-range-corrected $\anp$.

The $20\%$ zero-range residual is therefore not a ``CPP error'' at all; it is the distance any theory must travel from zero-range approximation to finite-range precision. What CPP \emph{has} supplied --- via $\Bd$ --- is the scale on which the whole construction rests.

Inverting the Bethe-Peierls relation $1/\anp = \kappa - \rzero \kappa^2/2 + \ldots$ using measured $\anp = 5.425$~fm and $\kappa = 0.2316$~fm$^{-1}$:
\[
\rzero = \frac{2(\kappa - 1/\anp)}{\kappa^2} = \frac{2(0.2316 - 0.1843)}{(0.2316)^2} = 1.76 \text{ fm}.
\]
The experimental effective range is $\rzero^{\mathrm{exp}} = 1.749$~fm \cite{effective_range}, so the leading-order effective-range expansion reproduces $\rzero$ to $+0.8\%$ using only $\Bd$ and $\anp$ as inputs. This is a notable consistency check: the deuteron is loosely enough bound that the effective-range expansion converges well at leading order, suggesting the underlying potential $\VSR(r)$ is smooth and reasonably characterized by the single length scale $\rzero$. Whatever CPP eventually derives for $\VSR(r)$ must be consistent with this observation.

\textbf{Category B (bipyramid-via-$\VSR$).}

\subsection{Singlet virtual state}

The $^1S_0$ np scattering channel has a virtual state near threshold: $E_v = +66$~keV above threshold (standard value). Analogously to the triplet, it is governed by the singlet channel $\VSR^{(s)}(r)$ which in base-to-base CPP corresponds to the K$_3$ face under orientation misalignment (see the $pp$/$nn$ discussion in SS-5). Quantitative prediction again requires $\VSR^{(s)}(r)$. \textbf{Category B.}

% ===================================================================
\section{Three-Category Classification}
\label{sec:classification}
% ===================================================================

Table~\ref{tab:classification} summarizes the classification for all deuteron observables considered.

\begin{table}[H]
\centering
\caption{Classification of deuteron observables by CPP derivability.}
\label{tab:classification}
\small
\begin{tabular}{@{}p{0.26\textwidth}p{0.18\textwidth}p{0.18\textwidth}p{0.28\textwidth}@{}}
\toprule
\textbf{Observable} & \textbf{Measured} & \textbf{CPP category} & \textbf{Status} \\
\midrule
Binding energy $\Bd$ & $2.2246$ MeV & A (bipyramid) & \textbf{Derived} in SS-5 at $+5.3\%$ \\
Spin $J^P = 1^+$ & $1^+$ & A (bipyramid) & Derived in SS-5 \\
Isospin $I = 0$ & $0$ & A (bipyramid) & Derived in SS-5 \\
Scattering length $\anp$ & $5.425$ fm & B (via $\VSR$) & Zero-range: $4.32$ fm ($-20\%$) \\
Effective range $\rzero$ & $1.749$ fm & B (via $\VSR$) & Inverted from $\anp, \kappa$: $1.76$ fm ($+0.8\%$) \\
Singlet virtual state & $+66$ keV & B (via $\VSR^{(s)}$) & Requires singlet $\VSR$ \\
Quadrupole $\Qd$ & $+0.286$ fm$^2$ & C (orbital) & Bipyramid lab-frame $-0.022$ fm$^2$: wrong sign, 10$\times$ too small \\
Matter radius $r_d$ & $1.975$ fm & C (orbital) & Orbital-scale \\
Charge radius $r_c$ & $2.141$ fm & C (orbital) & Orbital-scale \\
D-state admix $\PD$ & $4.5$--$5.8\%$ & C (orbital) & Requires orbital WF \\
Magnetic moment $\mud$ & $0.857 \mu_N$ & C (orbital, via $\PD$) & Requires orbital WF \\
\bottomrule
\end{tabular}
\end{table}

\begin{itemize}
\item \textbf{Category A (bipyramid-geometric, 3 observables):} binding, $J^P$, $I$. All derived in SS-5.
\item \textbf{Category B (bipyramid-via-$\VSR$, 3 observables):} $\anp$, $\rzero$, singlet virtual state. Derivable once OPEN-SS-20 is solved.
\item \textbf{Category C (orbital-dominated, 5 observables):} $\Qd$, $r_d$, $r_c$, $\PD$, $\mud$. Derivable once OPEN-SS-21 is solved.
\end{itemize}

The classification is substantive: it tells a would-be critic that the bipyramid mechanism makes \emph{no} claim about $\Qd$, $r_d$, $\PD$, or $\mud$ without additional CPP input. The Q$_d$-sign finding (Finding~\ref{find:Qd_sign}) \emph{supports} this classification --- the bipyramid actively predicts the wrong sign for $\Qd$ if one were to apply it naively, which reveals that a different physical mechanism must dominate.

% ===================================================================
\section{Open Problems}
\label{sec:open}
% ===================================================================

\begin{openproblem}[OPEN-SS-20: Derive $\VSR(r)$ shape from CPP primitives]
\label{open:VSR}
Starting from the base-to-base K$_3$ face structure of SS-5, derive the shape of the short-range np potential $\VSR(r)$ as a function of inter-nucleon separation $r$. Constraints:
\begin{itemize}
\item $\VSR(0) = -\Bpair = -2.342$ MeV (LO, rigid contact).
\item $\VSR(r) \to 0$ as $r \to \infty$.
\item Range scale: $\ledge = 0.364$ fm (natural CPP scale).
\end{itemize}
Candidate functional forms: Coulombic $\VSR(r) \propto -\ledge/\sqrt{\ledge^2 + r^2}$; Yukawa-like $\VSR(r) \propto -\exp(-r/\ledge)$; smooth cutoff $\VSR(r) \propto -\ledge/(\ledge + r)$. Each gives different numerical predictions for $\rzero$. The correct form must be selected by a CPP-structural argument, not a fit. Once derived, resolution gives:
\begin{itemize}
\item $\anp$ beyond zero range,
\item $\rzero$ from first principles,
\item singlet-channel $\VSR^{(s)}(r)$ and virtual-state energy,
\item np phase shifts at low energy.
\end{itemize}
\end{openproblem}

\begin{openproblem}[OPEN-SS-21: Deuteron orbital wavefunction from CPP]
\label{open:orbital}
Derive the deuteron relative-motion wavefunction $\psi_{np}(r)$ in the CPP framework, connecting the bipyramid core at $r \lesssim 1$~fm to the orbital extension at $r \sim 2$--$4$~fm. This is a multi-scale problem: at short range the K$_3$ face dynamics dominate; at long range the wavefunction is governed by the attractive potential tail plus the kinetic-energy balance appropriate for a loosely-bound state. Once resolved, yields:
\begin{itemize}
\item $r_d$, $r_c$ from $\langle r_{np}^2 \rangle$,
\item $\PD$ from $|w(r)|^2$ integrated weight,
\item $\Qd$ from standard quadrupole integrals including both intrinsic bipyramid contribution and orbital D-wave contribution,
\item $\mud$ via the $\PD$-dependent magnetic-moment relation.
\end{itemize}
The central unresolved question is how the CPP K$_3$ collective mode extends from the rigid-contact limit to finite inter-nucleon separation in a way compatible with the observed orbital size.
\end{openproblem}

% ===================================================================
\section{Discussion}
\label{sec:discussion}
% ===================================================================

\subsection{Why the negative finding is useful}

The Q$_d$-sign result (Finding~\ref{find:Qd_sign}) is not a failure of CPP; it is a clarification. Three points:

\textbf{(i)} It sharpens the mechanism's claims. SS-5 establishes that the bipyramid predicts binding; SS-6 establishes that the bipyramid does \emph{not} predict quadrupole. This is a falsifiable, testable distinction that any competing mechanism also has to navigate --- if an alternative theory claims to predict both $\Bd$ and $\Qd$ from the same rigid geometric object, that theory is suspect.

\textbf{(ii)} It matches conventional nuclear physics. In the deuteron, $\Qd$ is well-known to arise from the D-wave orbital tail, not from a short-range ``shape.'' CPP joining this understanding, rather than claiming to predict $\Qd$ from the bipyramid, is consistent with known physics.

\textbf{(iii)} It redirects programme effort. Instead of trying to squeeze $\Qd$, $r_d$, and $\PD$ out of the bipyramid, the programme now has a sharply defined open problem (OPEN-SS-21) --- derive the orbital wavefunction --- which is the right target for the next serious theoretical effort.

\subsection{The zero-range scattering length as a consistency check}

The $\anp = 1/\kappa = 4.32$~fm zero-range prediction is not a CPP-specific result. Any theory that correctly predicts $\Bd$ automatically gets this zero-range prediction right at the same precision, via the Bethe-Peierls relation. What CPP must do is \emph{go beyond} zero range by supplying $\VSR(r)$. The $20\%$ zero-range residual is therefore a measure of how far CPP has to travel to claim a precision prediction for $\anp$.

\subsection{Comparison with the SS-5 strategy}

SS-5's strength is a single mechanism giving seven concurrent independent predictions. SS-6's contribution is a single mechanism giving \emph{one} concurrent clarification: the bipyramid's active prediction of an oblate intrinsic $\Qd$ (wrong sign) reveals the orbital-dominance of the observed value. Compared to SS-5, this is a smaller paper with lower epistemic density, but it closes the easy avenue (``just predict more deuteron observables from the bipyramid'') cleanly and redirects effort toward the real open problems.

% ===================================================================
\section{Conclusion}
\label{sec:conclusion}
% ===================================================================

Applied to deuteron observables beyond the binding energy, the SS-5 base-to-base bipyramid mechanism yields three categories:
\begin{itemize}
\item \textbf{A (bipyramid-geometric):} binding, spin, isospin, parity --- derived in SS-5.
\item \textbf{B (bipyramid-via-$\VSR(r)$):} $\anp$, $\rzero$, singlet virtual state --- await $\VSR(r)$ derivation (OPEN-SS-20).
\item \textbf{C (orbital-dominated):} $\Qd$, $r_d$, $\PD$, $\mud$ --- await orbital wavefunction (OPEN-SS-21).
\end{itemize}

The rigid-bipyramid body-frame quadrupole is oblate, and its lab-frame contribution ($-0.022$ fm$^2$) is ten times smaller and opposite in sign to the measured $\Qd = +0.286$ fm$^2$ --- a diagnostic revealing that the observed $\Qd$ is orbital-dominated. Two universal scattering-theory consistency checks from $\Bd$ alone: zero-range $\anp = 1/\kappa = 4.32$ fm ($-20\%$ vs 5.43 fm), and effective-range inversion $\rzero = 2(\kappa - 1/\anp)/\kappa^2 = 1.76$ fm ($+0.8\%$ vs 1.749 fm). The 0.8\% $\rzero$ match is not a CPP-specific prediction but a notable consistency indicating the deuteron's short-range potential is smooth and well-characterized by a single length scale.

This paper does not close either of its two new open problems. Its contribution is scope clarification: it maps out which deuteron observables are downstream of the bipyramid (SS-5's territory), which require one further CPP input (OPEN-SS-20), and which require a new CPP framework for the orbital regime (OPEN-SS-21). The programme gains two sharp open problems and one definitive negative finding in exchange for an honest assessment of what the base-to-base mechanism actually claims.

\subsection{Problem status after this paper}

\begin{itemize}
\item OPEN-SS-10 (nuclear $V(r)$): clarified --- integrated binding at A=2,3,4 done in SS-5; short-range shape now separately OPEN-SS-20.
\item NEW OPEN-SS-20 ($\VSR(r)$ shape): OPEN.
\item NEW OPEN-SS-21 (orbital wavefunction): OPEN.
\item NEW PROP-SS-6-1 ($\Qd$ orbital-dominated; bipyramid lab-frame contribution 10$\times$ smaller and wrong sign): SUPPORTED by explicit body-frame and lab-frame calculations.
\item NEW PROP-SS-6-2 (universal Bethe-Peierls $\anp = 1/\kappa$ plus effective-range inversion): DERIVED from standard two-body theory. Zero-range $\anp = 4.32$~fm ($-20\%$); inverted $\rzero = 1.76$~fm ($+0.8\%$).
\end{itemize}

% ===================================================================
\section*{Acknowledgements}
% ===================================================================

The ``more stars for better navigation'' framing motivating this paper is Thomas Lee Abshier ND's (17 April 2026). The decision to write SS-6 as an honest scoping paper rather than claim derivations from orbital-dominated observables that the bipyramid does not predict is also Thomas's (18 April 2026). The Q$_d$-sign calculation and three-category classification were developed in Claude Opus's scoping session (18 April 2026) before the paper draft was begun; the scope was reduced from the initial pitch (P$_D$ + $\mud$ + $\anp$ + $\rzero$ as new CPP predictions) to the present honest classification after numerical exploration revealed that the orbital regime dominates four of the five targets.

% ===================================================================
\begin{thebibliography}{10}

\bibitem{abshier2026ss5}
T.~L.~Abshier, Claude Opus, ``Light-Nuclei Binding Energies from the Open-Vertex Cascade,'' SS-5 v6, Hyperphysics Institute (2026).

\bibitem{abshier2026ss2}
T.~L.~Abshier, Claude Opus, ``Lattice-Scale Grounding and Nucleon Structure from 600-Cell Geometry,'' SS-2 v1.0, Hyperphysics Institute (2026).

\bibitem{abshier2026sm8}
T.~L.~Abshier, Claude Opus, Grok, Copilot, ``Quark Generation Structure from 600-Cell Distance Shells,'' SM-8 v4.1, Hyperphysics Institute (2026).

\bibitem{pdg2024}
R.~L.~Workman et al.\ (Particle Data Group), Prog.\ Theor.\ Exp.\ Phys.\ \textbf{2022}, 083C01 (2022) and 2024 update.

\bibitem{ame2020}
M.~Wang et al., ``The AME 2020 atomic mass evaluation,'' Chin.\ Phys.\ C \textbf{45}, 030003 (2021).

\bibitem{krane_nuclear}
K.~S.~Krane, \emph{Introductory Nuclear Physics} (Wiley, 1987).

\bibitem{argonne_v18}
R.~B.~Wiringa, V.~G.~J.~Stoks, R.~Schiavilla, ``Accurate nucleon-nucleon potential with charge-independence breaking,'' Phys.\ Rev.\ C \textbf{51}, 38 (1995).

\bibitem{machleidt_entem}
R.~Machleidt, D.~R.~Entem, ``Chiral effective field theory and nuclear forces,'' Phys.\ Rep.\ \textbf{503}, 1 (2011).

\bibitem{anp_reference}
L.~Koester, W.~Nistler, ``New determination of the neutron-proton scattering lengths,'' Z.\ Phys.\ A \textbf{272}, 189 (1975).

\bibitem{effective_range}
S.~Klarsfeld, J.~Martorell, D.~W.~L.~Sprung, ``Deuteron properties and the nucleon-nucleon interaction,'' J.\ Phys.\ G \textbf{10}, 165 (1984).

\end{thebibliography}

% BEGIN GENERATED IDENTIFIER APPENDIX -- do not edit by hand
\clearpage
\section*{Appendix: Programme identifiers used in this paper}
\addcontentsline{toc}{section}{Appendix: Programme identifiers}

\noindent This programme tracks its results, assumptions and unresolved questions under short identifiers, so that a claim and its dependencies can be cited precisely. Prefixes denote the kind of item: \texttt{THEO-} a proved theorem, \texttt{FI-} a foundational input the derivation assumes, \texttt{PRED-} a registered prediction, \texttt{CONJ-} a conjecture, \texttt{OPEN-} an unresolved question, \texttt{PH-} a problem history. Those appearing in this paper are glossed below; no external document is needed to read them.

\begin{description}
  \item[\texttt{OPEN-SS-10}] \hfill \\ Derive nucleon–nucleon potential V(r) from qDP chain insertion dynamics.
  \item[\texttt{OPEN-SS-20}] \hfill \\ Derive the shape of the short-range nucleon-nucleon potential \$V\_\{\textbackslash\{\}mathrm\{SR\}\}(r)\$ as a function of inter-nucleon separation, starting from the base-to-base K\$\_3\$ face structure of SS-5.
  \item[\texttt{OPEN-SS-21}] \hfill \\ Derive the deuteron relative-motion wavefunction \$\textbackslash\{\}psi\_\{np\}(r)\$ connecting the bipyramid core at \$r \textbackslash\{\}lesssim 1\$ fm to the orbital extension at \$r \textbackslash\{\}sim 2\$-\$4\$ fm.
\end{description}
% END GENERATED IDENTIFIER APPENDIX

\end{document}
